\begin{lemma}{4.17}\label{III.4.17}
{\renewcommand{\thefootnote}{(109)}\nde{\normalfont One has modified 4.17 and~4.18, taking into account the additions made in~4.16.}}
Let \(X\) be an \(S\)-group smooth over~\(S\), and \(Y_{\mathcal{J}}\) a
sub-\(S_{\mathcal{J}}\)-group of \(X_{\mathcal{J}}\), flat and locally of finite
presentation over \(S_{\mathcal{J}}\). For each subscheme \(Y\) of~\(X\), flat
over~\(S\) and lifting \(Y_{\mathcal{J}}\),
% typo: missing space after the comma in ``relevant \(Y_{\mathcal{J}}\),consid\'erons''
let us consider the obstruction defined in~4.16.1:
\[
DY\in\Hom_{S_0}(Y_0\times_{S_0}Y_0,\,N_0)=C^{2}(Y_0,N_0)
\]
(i)~In order that \(Y\) be a sub-\(S\)-group of~\(X\), it is necessary and
sufficient that \(DY=0\).

(ii)~If one lets \(Y_0\) act on \(N_0\) by functoriality, starting from the
inner automorphisms of \(Y_0\), then \(DY\in Z^{2}(Y_0,N_0)\).

(iii)~If \(Y\) and \(Y'\) are two subschemes of~\(X\), flat over~\(S\), lifting
\(Y_{\mathcal{J}}\) (so that the deviation \(d(Y,Y')\in C^{1}(Y_0,N_0)\) is
defined, \cf~4.5.1), one has \(DY'-DY=\partial^{1}d(Y,Y')\).%
{\renewcommand{\thefootnote}{(110)}\nde{\normalfont In the original, one finds \(DY'-DY=-\partial d(Y,Y')\), but \(\partial\) there is the opposite of the differential \(\partial^{1}\) defined in~I, 5.1.}}
\end{lemma}

Let us prove these various assertions in succession.

\numpara{4.18}\label{III.4.18}
{\itshape Proof of 4.17~(i).}
By definition, one has \(DY=0\) if and only if \(Y\) is stable under the group
law of~\(X\). Thus \(DY=0\) if \(Y\) is a subgroup of~\(X\). Conversely, if
\(DY=0\), \(Y\) is equipped with the induced law \(\mu^{Y}\), which is
associative and reduces modulo \(\mathcal{J}\) according to the group law
on \(X_{\mathcal{J}}\); since \(Y\) is flat and locally of finite presentation
over~\(S\), it follows from~3.3 that \(\mu^{Y}\) is a group law.

\numpara{4.19}\label{III.4.19}
{\itshape Proof of 4.17~(ii).}
This is done by comparing the two values of
\(u=c(X,Y,\mu^{2}\circ(i,i,i))\) computed in the following two diagrams
\((D_{j})\), \(j=1,2\):
\[
(D_{j})\qquad
\xymatrix@C=2.6em@R=1.7em{
Y\times_{S}Y\times_{S}Y
  \ar@{^{(}->}[d]^{{(i,i,i)}}
&
Y\times_{S}Y
  \ar@{^{(}->}[d]^{{(i,i)}}
&
Y
  \ar@{^{(}->}[d]^{i}
\\
X\times_{S}X\times_{S}X
  \ar[r]_{f_{j}}
&
X\times_{S}X
  \ar[r]^{\mu}
&
X
}
\]
where \(f_{1}=(1,\pi)\), \(f_{2}=(\pi,1)\), and where one denotes by \(\mu^{2}\)
the morphism
\[
\mu\circ f_{1}=\mu\circ f_{2}\colon
X\times_{S}X\times_{S}X\longrightarrow X.
\]
% typo?: \(f_1=(1,\pi)\) and \(f_2=(\pi,1)\) here; below, \(f_1=(1,\mu)\) and \(f_2=(\mu,1)\). Kept as printed.
{\renewcommand{\thefootnote}{(111)}\nde{One has slightly modified the notations, and spelled out the beginning of the argument.}}
Set \(\mu^{Y}=\mu\circ(i,i)\), \(f_{j}^{Y}=f_{j}\circ(i,i,i)\) and
\(\mu^{2,Y}=\mu^{2}\circ(i,i,i)\). For \(j=1,2\), let us denote by \(a_{j}\)
and \(b_{j}\) the morphisms
\[
a_{j}=a_{\mu_{0}}\bigl((f_{j}^{Y})_{0}\bigr)
\qquad\text{and}\qquad
b_{j}=b_{(f_{j}^{Y})_{0}}\,,
\]
associated with the pair of morphisms \((f_{j}^{Y},\mu)\) by Lemma~4.12; one
therefore has:
\[
(\dagger)\qquad
\begin{cases}
\Hom_{\OO_{Y_{0}^{3}}}\bigl(
  (f_{j}^{Y})_{0}^{*}(\mathcal{N}_{Y_{0}\times Y_{0}/X_{0}\times X_{0}}),\,
  \mathcal{J}\,\OO_{Y_{0}^{3}}\bigr)
\xrightarrow{a_{j}}
\Hom_{\OO_{Y_{0}^{3}}}\bigl(
  (\mu^{2,Y})_{0}^{*}(\mathcal{N}_{Y_{0}/X_{0}}),\,
  \mathcal{J}\,\OO_{Y_{0}^{3}}\bigr)
\\[0.9em]
\Hom_{\OO_{Y_{0}^{2}}}\bigl(
  (\mu^{Y})_{0}^{*}(\mathcal{N}_{Y_{0}/X_{0}}),\,
  \mathcal{J}\,\OO_{Y_{0}^{2}}\bigr)
\xrightarrow{b_{j}}
\Hom_{\OO_{Y_{0}^{3}}}\bigl(
  (\mu^{2,Y})_{0}^{*}(\mathcal{N}_{Y_{0}/X_{0}}),\,
  \mathcal{J}\,\OO_{Y_{0}^{3}}\bigr)\,.
\end{cases}
\]

Since \(\mathcal{N}_{Y_{0}\times Y_{0}/X_{0}\times X_{0}}
\simeq\mathrm{pr}_{1}^{*}\mathcal{N}_{Y_{0}/X_{0}}
\oplus\mathrm{pr}_{2}^{*}\mathcal{N}_{Y_{0}/X_{0}}\)
(since \(X_{0}\) and \(Y_{0}\) are \emph{flat} over \(S_{0}\)), and
\(\mathcal{N}_{Y_{0}/X_{0}}\simeq\mathfrak{n}_{Y_{0}/X_{0}}\otimes_{\OO_{S_{0}}}\OO_{Y_{0}}\),
one has:
\[
(f_{j}^{Y})_{0}^{*}(\mathcal{N}_{Y_{0}\times Y_{0}/X_{0}\times X_{0}})
\simeq
(\mathfrak{n}_{Y_{0}/X_{0}}\oplus\mathfrak{n}_{Y_{0}/X_{0}})\otimes\OO_{Y_{0}^{3}}
\]
and, likewise,
\[
(\mu^{2,Y})_{0}^{*}(\mathcal{N}_{Y_{0}/X_{0}})
\simeq
\mathfrak{n}_{Y_{0}/X_{0}}\otimes\OO_{Y_{0}^{3}}
\qquad\text{and}\qquad
(\mu^{Y})_{0}^{*}(\mathcal{N}_{Y_{0}/X_{0}})
\simeq
\mathfrak{n}_{Y_{0}/X_{0}}\otimes\OO_{Y_{0}^{2}}\,.
\]
Moreover, since \(Y_{0}^{2}\) and \(Y_{0}^{3}\) are \emph{flat} over \(S_{0}\),
\((\dagger)\) is rewritten in the following form:
\[
(\ddagger)\qquad
\begin{cases}
a_{j}\colon\Hom_{S_{0}}(Y_{0}^{3},\,N_{0}\oplus N_{0})
  \to\Hom_{S_{0}}(Y_{0}^{3},\,N_{0})
\\
b_{j}\colon\Hom_{S_{0}}(Y_{0}^{2},\,N_{0})
  \to\Hom_{S_{0}}(Y_{0}^{3},\,N_{0})\,.
\end{cases}
\]
Applying~4.13 twice to \(c(X,Y,\mu^{2,Y})=u\), one obtains:
\[
a_{1}\bigl(c(X^{2},Y^{2},f_{1})\bigr)+b_{1}\bigl(c(X,Y,\mu^{Y})\bigr)
=
u
=
a_{2}\bigl(c(X^{2},Y^{2},f_{2})\bigr)+b_{2}\bigl(c(X,Y,\mu^{Y})\bigr)\,.
\]
Now, \(c(X,Y,\mu^{Y})=DY\) and, since \(f_{1}=(1,\mu)\) and \(f_{2}=(\mu,1)\),
one has, with evident notations:
\[
c(X^{2},Y^{2},f_{1})=(0,DY)
\qquad\text{and}\qquad
c(X^{2},Y^{2},f_{2})=(DY,0)\,.
\]
Thus one obtains:
\[
u=a_{1}((0,DY))+b_{1}(DY)=a_{2}((DY,0))+b_{2}(DY)\,.
\]
The first thing that one notices is that \(b_{j}\) is none other than
\(\Hom_{S_{0}}((f_{j}^{Y})_{0},N_{0})\), that is, the morphism deduced from
\((f_{j}^{Y})_{0}\) by functoriality.

The identity above therefore becomes:
\[
a_{1}((0,DY))-\Hom((\mu,1),N_{0})(DY)
+\Hom((1,\mu),N_{0})(DY)
-a_{2}((DY,0))=0\,.
\]
One recognizes the two middle terms: they are the parts ``\(DY(xy,z)\)'' and
``\(DY(x,yz)\)'' of the formula of the \(2\)-coboundary. It therefore only
remains to identify the two other terms.

We must first compute the map \(a_{j}\). Now it arises, by inverse image
under \((f_{j}^{Y})_{0}\), from the morphism of \(\OO_{Y_{0}^{2}}\)-modules
\[
P\colon
\mathfrak{n}_{Y_{0}/X_{0}}\otimes\OO_{Y_{0}^{2}}
\longrightarrow
(\mathfrak{n}_{Y_{0}/X_{0}}\oplus\mathfrak{n}_{Y_{0}/X_{0}})\otimes\OO_{Y_{0}^{2}}
\]
induced by the product in \(Y_{0}\). Now this morphism is described in the
following way: let us consider the vector bundle
\(\mathbf{V}=\mathbf{V}(\mathfrak{n}_{Y_{0}/X_{0}})\); \(P\) gives by duality a
morphism
\[
\mathbf{V}(P)\colon
\mathbf{V}\times_{S_{0}}\mathbf{V}\times_{S_{0}}Y_{0}\times_{S_{0}}Y_{0}
\longrightarrow
\mathbf{V}\times_{S_{0}}Y_{0}\times_{S_{0}}Y_{0}
\]
which is expressed set-theoretically by
\[
\mathbf{V}(P)(u,v,a,b)=(u+\Ad(a)v,\,ab,\,b)\,.
\]
{\renewcommand{\thefootnote}{(112)}\nde{One has replaced \(a\), \(b\) by \(ab\), \(b\), in order to make it plain that \(\mathbf{V}(P)\) arises by inverse image on \(Y_{0}^{2}\) from the multiplication morphism \(\mathbf{V}_{Y_{0}}\times_{S_{0}}\mathbf{V}_{Y_{0}}\to\mathbf{V}_{Y_{0}}\).}}
This is proved exactly like the corresponding fact for Lie algebras, that is,
for the module \(\omega^{1}_{Y_{0}/S_{0}}\). One first observes that
\(\mathbf{V}\) is equipped, by functoriality in \(Y_{0}\), with a group
structure in the category of vector bundles over \(S_{0}\); by virtue of the
lemma already used for Lie algebras (Expos\'e~II, 3.10), this structure
coincides with the group structure underlying its structure of
\(\OO_{S}\)-module.
% typo?: printed \(\OO_S\)-module; the surrounding base is \(S_0\). Kept as printed.
One then sees that
\(\mathbf{V}(\mathfrak{n}_{Y_{0}/X_{0}}\otimes_{\OO_{S_{0}}}\OO_{Y_{0}})
=\mathbf{V}(\mathcal{N}_{Y_{0}/X_{0}})\)
is likewise equipped with a structure of \(S_{0}\)-group which is none other
than the semidirect product of that of \(\mathbf{V}\) by that of \(Y_{0}\).
It only remains to identify the operations of \(Y_{0}\) on \(\mathbf{V}\) in
order to establish the desired formula.

Let us now compute the two remaining terms. Let us first consider
\(a_{1}((0,DY))\). One computes it by the diagram (where \(\mathfrak{n}\) denotes
\(\mathfrak{n}_{Y_{0}/X_{0}}\)):
\[
\xymatrix@C=5.5em@R=1.55em{
\mathfrak{n}\otimes\OO_{Y_{0}^{2}}
  \ar[r]^{P}
  \ar[dd]_{(f_{1}^{Y})_{0}^{*}}
&
(\mathfrak{n}+\mathfrak{n})\otimes\OO_{Y_{0}^{2}}
  \ar[dd]^{(f_{1}^{Y})_{0}^{*}}
\\
\\
\mathfrak{n}\otimes\OO_{Y_{0}^{3}}
  \ar[dr]_{{a_{1}((0,DY))}}
&
(\mathfrak{n}+\mathfrak{n})\otimes\OO_{Y_{0}^{3}}
  \ar[d]^{{(0,DY)}}
\\
&
\mathcal{J}\otimes\OO_{Y_{0}^{3}}\,.
}
\]
Considering now the vector bundles defined by these various modules as so
many schemes over \(S_{0}\), and taking the points with values in anything
whatsoever, one has, denoting by \((u,x,y,z)\) a point of
\(\mathbf{V}(\mathcal{J})\times Y_{0}^{3}\):
\[
\xymatrix@C=7em@R=2.5em{
\bigl(\Ad(x)\,DY_{y,z}(u),\,x,\,yz\bigr)
&
\bigl(0+DY_{y,z}(u),\,x,\,yz\bigr)
  \ar[l]
\\
&
\bigl(0+DY_{y,z}(u),\,x,\,y,\,z\bigr)
  \ar[u]
\\
\bigl(\Ad(x)\,DY_{y,z}(u),\,x,\,y,\,z\bigr)
  \ar[uu]
&
\bigl(u,\,x,\,y,\,z\bigr)
  \ar[l]
  \ar[u]
}
\]
One has therefore obtained
\(a_{1}((0,DY))(x,y,z)=\Ad(x)\,DY(y,z)\), which is indeed the first term of
the coboundary. One would likewise have
\(a_{2}((DY,0))(x,y,z)=DY(x,y)\), whence%
{\renewcommand{\thefootnote}{(113)}\nde{One has changed the signs, in order to make them compatible with~I 5.1.}}
\[
0=\Ad(x)\,DY(y,z)-DY(xy,z)+DY(x,yz)-DY(x,y)
=(\partial^{2}DY)(x,y,z)\,.
\]

\numpara{4.20}\label{III.4.20}
{\itshape Proof of 4.17~(iii).}%
{\renewcommand{\thefootnote}{(114)}\nde{One has spelled out the original in what follows.}}
This is done by comparing the two values of
\(v=c(X,Y,\mu\circ(i',i'))\) computed in the following two diagrams
\[
(*)\qquad
\xymatrix@C=3.2em@R=1.6em{
&
Y'
  \ar@{^{(}->}[d]^{i'}
&
Y
  \ar@{^{(}->}[d]^{i}
\\
Y'\times_{S}Y'
  \ar[r]^{{\mu\circ(i',i')}}
&
X
&
X
}
\]
\[
(\dagger)\qquad
\xymatrix@C=3.2em@R=1.6em{
&
Y\times_{S}Y
  \ar@{^{(}->}[d]^{{(i,i)}}
&
Y
  \ar@{^{(}->}[d]^{i}
\\
Y'\times_{S}Y'
  \ar[r]^{{(i',i')}}
&
X\times_{S}X
  \ar[r]^{\mu}
&
X\,.
}
\]
Let us denote \(f=\mu\circ(i',i')\); then \((*)\) gives
\begin{equation*}
(1)\qquad
v=DY'+f_{0}^{*}\bigl(c(X,Y,i')\bigr)\,.
\end{equation*}
Now \(Y'_{0}=Y_{0}\) and \(f_{0}\) is the multiplication \(Y_{0}^{2}\to Y_{0}\);
one deduces from this that
\begin{equation*}
(2)\qquad
f_{0}^{*}\bigl(c(X,Y,i')\bigr)(x_{0},y_{0})
=
c(X,Y,i')(x_{0}y_{0})\,.
\end{equation*}
Set \(c=c(X,Y,i')\); via the identification \(N'_{0}\simeq N_{0}\oplus N_{0}\),
\(c(X\times_{S}X,\,Y\times_{S}Y,\,(i',i'))\) is identified with the pair
\((c,c)\). Then, denoting \(h=(i',i')\), \((\dagger)\) gives
\begin{equation*}
(3)\qquad
v=h_{0}^{*}(DY)+a_{\mu_{0}}(c,c)\,.
\end{equation*}
Now \(h_{0}\) is the identity map of \(Y_{0}^{2}\), whence
\(h_{0}^{*}(DY)=DY\). Finally, by the computation of \(a_{\mu_{0}}\) carried
out previously, one has, for every \(S'\to S\) and
\(x_{0},\,y_{0}\in Y_{0}(S'_{0})\),
\begin{equation*}
(4)\qquad
a_{\mu_{0}}(c,c)(x_{0},y_{0})
=
c(x_{0})+\Ad(x_{0})\bigl(c(y_{0})\bigr)\,.
\end{equation*}
One therefore obtains:
\begin{align*}
(DY'-DY)(x_{0},y_{0})
&=
\Ad(x_{0})\bigl(c(X,Y,i')(y_{0})\bigr)
-c(X,Y,i')(x_{0}y_{0})
+c(X,Y,i')(x_{0})
\\
&=
\bigl(\partial^{1}c(X,Y,i')\bigr)(x_{0},y_{0})\,.
\end{align*}
Since \(c(X,Y,i')=d(Y,Y')\) (\cf~4.8~(ii)), this shows that
\(DY'-DY=\partial^{1}d(Y,Y')\).
% typo: printed \(\partial^{1}d((Y,Y')\) -> \(\partial^{1}d(Y,Y')\)

\begin{theorem}{4.21}\label{III.4.21}
Let \(S\) be a scheme, \(\mathcal{I}\) and \(\mathcal{J}\) two ideals%
{\renewcommand{\thefootnote}{(115)}\nde{\normalfont One has removed the hypothesis that \(\mathcal{I}\) be nilpotent, which seems superfluous.}}
on \(S\) such that \(\mathcal{I}\supset\mathcal{J}\) and
\(\mathcal{I}\cdot\mathcal{J}=0\). Let \(X\) be an \(S\)-group smooth over~\(S\),
and \(Y_{\mathcal{J}}\) a sub-\(S_{\mathcal{J}}\)-group of \(X_{\mathcal{J}}\),
flat and locally of finite presentation over \(S_{\mathcal{J}}\). Let us
consider the commutative \(S_{0}\)-group functor \(N_{0}\) defined by
\[
\Hom_{S_{0}}(T,N_{0})
=
\Hom_{\OO_{T}}(\mathfrak{n}_{Y_{0}/X_{0}}\otimes_{\OO_{S_{0}}}\OO_{T},\,
\mathcal{J}\otimes_{\OO_{S_{0}}}\OO_{T})\,,
\qquad
T\in\Ob\Sch_{/S_{0}}\,,
\]
on which \(Y_{0}\) acts by means of the inner automorphisms of \(X_{0}\).

(i)~In order that there exist a sub-\(S\)-group of~\(X\), flat over~\(S\),
which reduces to \(Y_{\mathcal{J}}\), it is necessary and sufficient that the
following two conditions be satisfied:

(i\(_1\))~There exists a subscheme \(Y\) of~\(X\), flat over~\(S\), lifting
\(Y_{\mathcal{J}}\) (a condition automatically satisfied if \(Y_{0}\) is
affine, \cf~4.6.5).

(i\(_2\))~A certain canonical obstruction, an element of
\(H^{2}(Y_{0},N_{0})\), is zero.

(ii)~If the conditions of~(i) are satisfied, the set of sub-\(S\)-groups \(Y\)
of~\(X\), flat over~\(S\) and reducing to \(Y_{\mathcal{J}}\), is a principal
homogeneous set under the group \(Z^{1}(Y_{0},N_{0})\).%
{\renewcommand{\thefootnote}{(116)}\nde{\normalfont The question whether the preceding set, modulo conjugation by the \(x\in X(S)\) inducing the unit element of \(X(S_{\mathcal{J}})\), is principal homogeneous under \(H^{1}(Y_{0},N_{0})\), occupies nos.~4.23 to~4.36.}}
\end{theorem}

Indeed, condition (i\(_1\)) is necessary. Let us assume it satisfied, and let
\(Y\) be flat over~\(S\) and lifting \(Y_{\mathcal{J}}\). We must seek a \(Y'\)
flat over~\(S\), also lifting \(Y_{\mathcal{J}}\), such that \(DY'=0\),%
{\renewcommand{\thefootnote}{(117)}\nde{One has corrected \(\partial DY'\) to \(DY'\).}}
\cf~4.17~(i). By~4.17~(iii), this amounts to seeking a
\(d(Y',Y)\in C^{1}(Y_{0},N_{0})\) such that
\(DY=\partial^{1}d(Y',Y)\).%
{\renewcommand{\thefootnote}{(118)}\nde{\normalfont \Cf\ N.D.E.~(110).}}

Let \(c\in H^{2}(Y_{0},N_{0})\) be the image class of \(DY\), which is a
cocycle by~4.17~(ii). It does not depend on the choice of \(Y\) by~4.17~(iii),
and its vanishing is necessary and sufficient for the existence of a
\(d(Y',Y)\) satisfying the preceding equation. This proves~(i).

If one has now chosen \(Y\) such that \(DY=0\), the equation to be solved
reads \(\partial^{1}d(Y',Y)=0\), which proves~(ii).

\begin{remark}{4.22}\label{III.4.22}
One keeps the notations of~4.21. By~4.15, \(Y_{0}\) is locally a complete
intersection in \(X_{0}\), so \(\mathcal{N}_{Y_{0}/X_{0}}\) is a finite
locally free \(\OO_{Y_{0}}\)-module, and consequently
\(\mathfrak{n}_{Y_{0}/X_{0}}=\varepsilon_{0}^{*}(\mathcal{N}_{Y_{0}/X_{0}})\)
is a finite locally free \(\OO_{S_{0}}\)-module. Thus, denoting
\(\mathfrak{n}^{\vee}_{Y_{0}/X_{0}}
=\mathcal{H}\mathrm{om}_{\OO_{Y_{0}}}(\mathfrak{n}_{Y_{0}/X_{0}},\OO_{Y_{0}})\),
% typo?: dual written as \(\mathcal{H}\mathrm{om}\) over \(\OO_{Y_0}\). Kept as printed.
one has
\[
\mathcal{H}\mathrm{om}_{\OO_{T}}
  (\mathfrak{n}_{Y_{0}/X_{0}}\otimes_{\OO_{S_{0}}}\OO_{T},\,
   \mathcal{J}\otimes_{\OO_{S_{0}}}\OO_{T})
\simeq
\mathfrak{n}^{\vee}_{Y_{0}/X_{0}}\otimes_{\OO_{S_{0}}}\mathcal{J}
  \otimes_{\OO_{S_{0}}}\OO_{T}
\]
for every \(T\to S_{0}\).%
{\renewcommand{\thefootnote}{(119)}\nde{One has spelled out what precedes; this shows that the isomorphism which follows is valid without a flatness hypothesis. On the other hand, since~4.16, one has restricted oneself to the \(S\)-schemes \(f\colon T\to S\) \emph{flat} over~\(S\), in order to ensure that the group \(\Hom_{\OO_{T_{0}}}(\mathfrak{n}_{Y_{0}/X_{0}}\otimes_{\OO_{S_{0}}}\OO_{T_{0}},\,\mathcal{J}\,\OO_{T})\), in which the obstruction \(c(X,Y,f)\) lies, coincides with \(N_{0}(T_{0})\) (\cf\ the end of~4.16).}}
Consequently, the \(S_{0}\)-functor \(N_{0}\) is isomorphic to the functor
\[
\mathbf{W}(\mathfrak{n}^{\vee}_{Y_{0}/X_{0}}\otimes_{\OO_{S_{0}}}\mathcal{J})
\simeq
\mathbf{W}\bigl(\mathcal{H}\mathrm{om}_{\OO_{S_{0}}}
  (\mathfrak{n}_{Y_{0}/X_{0}},\mathcal{J})\bigr)\,.
\]
There result the isomorphisms:%
{\renewcommand{\thefootnote}{(120)}\nde{with the notations of~I 5.3, assuming \(Y_{0}\) affine over \(S_{0}\).}}
\begin{align*}
H^{2}(Y_{0},N_{0})
&\simeq
H^{2}\bigl(Y_{0},\,
  \mathcal{H}\mathrm{om}_{\OO_{S_{0}}}(\mathfrak{n}_{Y_{0}/X_{0}},\mathcal{J})\bigr)
\simeq
H^{2}(Y_{0},\,\mathfrak{n}^{\vee}_{Y_{0}/X_{0}}\otimes_{\OO_{S_{0}}}\mathcal{J})\,,
\\
Z^{1}(Y_{0},N_{0})
&\simeq
Z^{1}\bigl(Y_{0},\,
  \mathcal{H}\mathrm{om}_{\OO_{S_{0}}}(\mathfrak{n}_{Y_{0}/X_{0}},\mathcal{J})\bigr)
\simeq
Z^{1}(Y_{0},\,\mathfrak{n}^{\vee}_{Y_{0}/X_{0}}\otimes_{\OO_{S_{0}}}\mathcal{J})\,.
\end{align*}
% typo: an extra closing parenthesis after each of the two displays. Removed.
\end{remark}

\numpara{4.23}\label{III.4.23}
Still under the hypotheses of~4.21, we shall now study how the set of the
\(Y\) lifting \(Y_{\mathcal{J}}\) behaves with respect to conjugation by
sections of~\(X\). If \(x\) is a section of \(X\) over \(S\) inducing the unit
section of \(X_{\mathcal{J}}\), the inner automorphism \(\Int(x)\) defined
by~\(x\) transforms flat subgroups of \(X\) lifting \(Y_{\mathcal{J}}\) into
flat subgroups of \(X\) lifting \(Y_{\mathcal{J}}\). Now, under the conditions
of~4.21~(ii), the set of these subgroups is principal homogeneous under
\(Z^{1}(Y_{0},N_{0})\); we shall see that there then exists a subgroup
\(\Gamma\) of \(B^{1}(Y_{0},N_{0})\)%
{\renewcommand{\thefootnote}{(121)}\nde{One has replaced \(Z^{1}(Y_{0},N_{0})\) by \(B^{1}(Y_{0},N_{0})\), since the proof shows that \(\Gamma\) is a subgroup of \(B^{1}(Y_{0},N_{0})\), \cf~4.27--4.29.}}
such that two subgroups of~\(X\), flat over~\(S\) and lifting
\(Y_{\mathcal{J}}\), are conjugate (by \(x\in X(S)\) inducing the unit element
of \(X(S_{\mathcal{J}})\)) if and only if their ``difference'' in
\(Z^{1}(Y_{0},N_{0})\) is an element of~\(\Gamma\). In the most favorable
cases, we shall show that \(\Gamma\) equals \(B^{1}(Y_{0},N_{0})\), so that
the set of flat subgroups of \(X\) lifting \(Y_{\mathcal{J}}\), modulo
conjugation by the \(x\in X(S)\) inducing the unit element of
\(X(S_{\mathcal{J}})\), is empty or principal homogeneous under
\(H^{1}(Y_{0},N_{0})\) (\cf~4.29 and~4.36).

\numpara{4.24}\label{III.4.24}
One keeps the notations of~4.21. Let \(Y\) be a flat subgroup of~\(X\),
reducing to \(Y_{\mathcal{J}}\). Let us recall that we introduced in~0.5 the
functor \(L_{0X}\) (\resp\ \(L_{0Y}\)) defined by the following identity with
respect to the variable \(S_{0}\)-scheme~\(T\):
\[
\Hom_{S_{0}}(T,\,L_{0X})
=
\Hom_{\OO_{T}}(\omega^{1}_{X_{0}/S_{0}}\otimes_{\OO_{S_{0}}}\OO_{T},\,
\mathcal{J}\otimes_{\OO_{S_{0}}}\OO_{T})
\]
(\resp\ likewise upon replacing \(X\) by~\(Y\)), as well as the functor
\(L'_{X}=\prod_{S_{0}/S}L_{0X}\).

Now one has:

\begin{lemma}{4.25}\label{III.4.25}
There exists a canonical exact sequence of \(Y_{0}\)-\(\OO_{S_{0}}\)-modules
\[
(+)\qquad
\mathfrak{n}_{Y_{0}/X_{0}}
\xrightarrow{d}
\omega^{1}_{X_{0}/S_{0}}
\longrightarrow
\omega^{1}_{Y_{0}/S_{0}}
\longrightarrow
0
\]
having the following properties:

(i)~By inverse image on \(Y_{0}\), \(d\) gives the morphism \(D_{0}\) of
4.8~(iii).

(ii)~If \(X_{0}\) and \(Y_{0}\) are smooth over \(S_{0}\), then \(d\) is
injective. Since the two \(\omega^{1}\) are then locally free of finite type,
the same holds for \(\mathfrak{n}_{Y_{0}/X_{0}}\), and the sequence is locally
split.
\end{lemma}

\begin{proof}
{\renewcommand{\thefootnote}{(122)}\nde{One has spelled out the proof, taking into account the additions made in Exp.~I, \S~6.8.}}
Let \(\pi_{0}\) denote the morphism \(Y_{0}\to S_{0}\). By SGA~1~II,
formula~(4.3) (see also EGA~IV\(_4\), 16.4.21), one has a canonical exact
sequence of \(\OO_{Y_{0}}\)-modules
\[
(\dagger)\qquad
\mathcal{N}_{Y_{0}/X_{0}}
\xrightarrow{D_{0}}
\Omega^{1}_{X_{0}/S_{0}}\otimes_{\OO_{X_{0}}}\OO_{Y_{0}}
\longrightarrow
\Omega^{1}_{Y_{0}/S_{0}}
\longrightarrow
0\,.
\]
Since this sequence is formed of \((Y_{0}\times_{S}Y_{0})\)-equivariant
modules and morphisms, its inverse image \((+)\) by \(\varepsilon_{0}^{*}\) is
an exact sequence of \(Y_{0}\)-\(\OO_{S_{0}}\)-modules, and \((\dagger)\) is the
inverse image of \((+)\) by \(\pi_{0}^{*}\) (\cf\ Exp.~I, \S~6.8). This
proves~(i).

Suppose moreover that \(X_{0}\) and \(Y_{0}\) are smooth over \(S_{0}\). Then,
by SGA~1~II~4.10 (see also EGA~IV\(_4\), 17.2.3~(i) and~17.2.5), \(D\) is
injective and the sequence \((\dagger)\) is formed of
\(\OO_{Y_{0}}\)-modules locally free of finite type (hence is locally split).
By the equivalence of categories I,~6.8.1, \(d\) is likewise injective, and
thus the sequence \((+)\) has the properties indicated.
% typo?: the arrow of \((\dagger)\) is labeled \(D_0\), while the prose says ``\(D\) is injective''. Kept as printed.
\end{proof}

\numpara{4.26}\label{III.4.26}
{\renewcommand{\thefootnote}{(123)}\nde{One has spelled out the original, in order to make it plain that one has an exact sequence of \(Y_{0}\)-\(\OO_{S_{0}}\)-modules.}}
For every \(S_{0}\)-scheme \(f\colon T\to S_{0}\), \((+)\) gives an exact
sequence of \(Y_{0}(T)\)-\(\OO(T)\)-modules
\[
\begin{aligned}
0&\longrightarrow
\Hom_{\OO_{T}}\bigl(f^{*}(\omega^{1}_{Y_{0}/S_{0}}),\,f^{*}(\mathcal{J})\bigr)
\longrightarrow
\Hom_{\OO_{T}}\bigl(f^{*}(\omega^{1}_{X_{0}/S_{0}}),\,f^{*}(\mathcal{J})\bigr)
\\
&\longrightarrow
\Hom_{\OO_{T}}\bigl(f^{*}(\mathfrak{n}_{Y_{0}/X_{0}}),\,f^{*}(\mathcal{J})\bigr)\,,
\end{aligned}
\]
one therefore has an exact sequence of \(Y_{0}\)-\(\OO_{S_{0}}\)-modules:
\begin{equation}
0\longrightarrow L_{0Y}\longrightarrow L_{0X}\xrightarrow{d}N_{0}\,.
\tag{4.26.1}
\end{equation}
One deduces from this an exact sequence of complexes of abelian groups:
\[
0\longrightarrow
C^{*}(Y_{0},L_{0Y})\longrightarrow
C^{*}(Y_{0},L_{0X})\xrightarrow{d^{*}}
C^{*}(Y_{0},N_{0})\,,
\]
and in particular, a commutative diagram with exact rows
\[
\xymatrix@C=1.5em@R=1.8em{
0\ar[r]
&
C^{0}(Y_{0},L_{0Y})
  \ar[r]
  \ar[d]^{\partial}
&
C^{0}(Y_{0},L_{0X})
  \ar[r]^{d^{0}}
  \ar[d]^{\partial}
&
C^{0}(Y_{0},N_{0})
  \ar[d]^{\partial}
\\
0\ar[r]
&
C^{1}(Y_{0},L_{0Y})
  \ar[r]
&
C^{1}(Y_{0},L_{0X})
  \ar[r]^{d^{1}}
&
C^{1}(Y_{0},N_{0})\,.
}
\]
Let us note that \(C^{0}(Y_{0},L_{0Y})\) (\resp\ \(C^{0}(Y_{0},L_{0X})\)) is
none other than \(\Hom_{S_{0}}(S_{0},L_{0Y})=\Hom_{S}(S,L'_{Y})\)
(\resp\ \(\cdots\)), \ie\ (\cf~0.9) the group of sections of \(Y\) (\resp\
of~\(X\)) over \(S\) inducing the unit section of \(X_{\mathcal{J}}\). Let us
also note that \(d^{1}\) is none other than the morphism \(v_{i_{Y_{0}}}\)
of~4.8~(iii), where \(i_{Y_{0}}\colon Y_{0}\to X_{0}\) is the canonical
immersion.%
{\renewcommand{\thefootnote}{(124)}\nde{This follows from the definition of \(d\colon\mathfrak{n}_{Y_{0}/X_{0}}\to\Omega^{1}_{X_{0}/S_{0}}\) (\cf~4.25) and from that of \(v_{i_{Y_{0}}}\) (\cf~4.8).}}
% typo?: footnote~(124) prints \(\Omega^{1}_{X_0/S_0}\); the map \(d\) of~4.25 lands in \(\omega^{1}_{X_0/S_0}\). Kept as printed.

\begin{lemma}{4.27}\label{III.4.27}
Under the conditions of~4.21 for \(S\), \(\mathcal{I}\), \(\mathcal{J}\)
and~\(X\), let \(Y\) be a subgroup of~\(X\), flat over~\(S\) and lifting
\(Y_{\mathcal{J}}\). Let \(i\colon Y\hookrightarrow X\) denote the canonical
immersion.%
{\renewcommand{\thefootnote}{(125)}\nde{\normalfont One has added point~(i) below, which will be useful in~4.35.1 and then in~4.38~(4) and~(5).}}

(i)~Let \(i'\colon Y\to X\) be a morphism of \(S\)-schemes lifting \(i_{0}\)
(so that \(i'\) is also an immersion), let \(Y'=i'(Y)\) and let \(d(i,i')\) be
the element of \(C^{1}(Y_{0},L_{0X})\) such that \(i'=d(i,i')\cdot i\)
(\cf~1.2.4). Then the deviation \(d(Y,Y')\in C^{1}(Y_{0},N_{0})\)
(\cf~4.5.1) is given by the formula:
\[
d(Y,Y')=d^{1}\bigl(d(i,i')\bigr)\,.
\]

(ii)~Let \(x\in C^{0}(Y_{0},L_{0X})\) be a section of \(X\) over \(S\) inducing
the unit section of \(X_{\mathcal{J}}\) over \(S_{\mathcal{J}}\). Then the
deviation \(d(Y,\Int(x)Y)\in C^{1}(Y_{0},N_{0})\) (\cf~4.5.1) is given by the
formula:
\[
-d\bigl(Y,\Int(x)Y\bigr)=d^{1}\partial x=\partial\,d^{0}x\,.
\]
\end{lemma}

Indeed, \(Y'\) is the image of \(Y\) by the composite morphism:%
{\renewcommand{\thefootnote}{(126)}\nde{One recalls that \(L'_{X}=\prod_{S_{0}/S}L_{0X}\).}}
\[
Y\xrightarrow{(d(i,i'),\,i)}L'_{X}\times_{S}X\longrightarrow X\,,
\]
which is denoted \(d(i,i')\cdot i\) in~4.8~(iii); by \loccit\ and the equality
\(v_{i_{0}}=d^{1}\), one therefore has:
\[
c\bigl(X,Y',\,d(i,i')\cdot i\bigr)-c(X,Y',i)
=
v_{i_{0}}\bigl(d(i,i')\bigr)
=
d^{1}\bigl(d(i,i')\bigr)\,.
\]
But \(d(i,i')\cdot i=i'\) factors through \(Y'\) by definition, so the first
term is zero; moreover, by~4.8~(ii), one has
\(c(X,Y',i)=d(Y',Y)=-d(Y,Y')\). Thus
\(d(Y,Y')=d^{1}\bigl(d(i,i')\bigr)\), which proves~(i).

Now let \(x\) be as in~(ii). By the formula
\[
xyx^{-1}=xyx^{-1}y^{-1}y=(x-\Ad(y)x)y=(-\partial x)(y)\cdot y\,,
\]
one sees that \(Y'\) is the image of \(Y\) by the immersion
\(i'=(-\partial x)\cdot i_{Y}\). Thus, by~(i), one obtains
\[
-d\bigl(Y,\Int(x)Y\bigr)=d^{1}\partial x=\partial\,d^{0}x\,.
\]

\begin{corollary}{4.28}\label{III.4.28}
In order that two subgroups \(Y\) and \(Y'\) of~\(G\), flat over~\(S\) and
lifting \(Y_{\mathcal{J}}\), be conjugate by a section of \(X\) over \(S\)
inducing the unit section of \(X_{\mathcal{J}}\), it is necessary and
sufficient that
\(d(Y,Y')\in\partial\,d^{0}\,C^{0}(Y_{0},L_{0X})
\subset\partial\,C^{0}(Y_{0},N_{0})=B^{1}(Y_{0},N_{0})\).
% typo?: printed ``of \(G\)''; the surrounding group is \(X\). Kept as printed.
\end{corollary}

\begin{corollary}{4.29}\label{III.4.29}
If \(d^{0}\) is surjective, \(Y\) and \(Y'\) as above are conjugate by a
section of \(X\) over \(S\) inducing the unit section of \(X_{\mathcal{J}}\)
if and only if \(d(Y,Y')\in B^{1}(Y_{0},N_{0})\).
\end{corollary}

\begin{corollary}{4.30}\label{III.4.30}
Let \(Y\) be as in~4.27; the set of conjugates of \(Y\) by sections of \(X\)
over \(S\) inducing the unit section of \(X_{\mathcal{J}}\) is isomorphic to:
\[
d^{1}\partial\,C^{0}(Y_{0},L_{0X})
=
C^{0}(Y_{0},L_{0X})\big/\Ker d^{1}\partial\,.
\]
\end{corollary}

Let us now note that \(C^{0}(Y_{0}/L_{0X})/\Ker d^{1}\partial\) is computed
solely by means of the left-hand square of the commutative diagram of~4.26.
% typo?: printed \(C^{0}(Y_{0}/L_{0X})\); the line above has \(C^{0}(Y_{0},L_{0X})\). Kept as printed.
It follows in particular that one can also compute it in any diagram of the
same type having the same left-hand square. Let us consider in particular the
functor \(L_{0X}/L_{0Y}\) over \(S_{0}\) defined by
\[
\Hom_{S_{0}}(T,\,L_{0X}/L_{0Y})
=
\Hom_{S_{0}}(T,\,L_{0X})/\Hom_{S_{0}}(T,\,L_{0Y})\,.
\]
One has a commutative diagram
\[
\xymatrix@C=1.35em@R=1.8em{
0\ar[r]
&
C^{0}(Y_{0},L_{0Y})
  \ar[r]
  \ar[d]^{\partial}
&
C^{0}(Y_{0},L_{0X})
  \ar[r]
  \ar[d]^{\partial}
&
C^{0}(Y_{0},L_{0X}/L_{0Y})
  \ar[r]
  \ar[d]^{\partial}
&
0
\\
0\ar[r]
&
C^{1}(Y_{0},L_{0Y})
  \ar[r]
&
C^{1}(Y_{0},L_{0X})
  \ar[r]
&
C^{1}(Y_{0},L_{0X}/L_{0Y})
  \ar[r]
&
0\,,
}
\]
whence, by the preceding remark:

\begin{corollary}{4.31}\label{III.4.31}
Let \(Y\) be as in~4.27; the set of conjugates of \(Y\) by sections of \(X\)
over \(S\) inducing the unit section of \(X_{\mathcal{J}}\) is isomorphic to
\[
E=\partial\,C^{0}(Y_{0},\,L_{0X}/L_{0Y})
=
C^{0}(Y_{0},\,L_{0X}/L_{0Y})
\big/
H^{0}(Y_{0},\,L_{0X}/L_{0Y})\,.
\]
\end{corollary}

\begin{corollary}{4.32}\label{III.4.32}
Suppose moreover that \(S_{0}\) is affine and that \(\omega^{1}_{Y_{0}/S_{0}}\)%
{\renewcommand{\thefootnote}{(127)}\nde{\normalfont One has corrected \(\mathcal{L}\mathrm{ie}(Y_{0}/S_{0})\) to \(\omega^{1}_{Y_{0}/S_{0}}\).}}
is finite locally free. If one denotes
\(\mathcal{F}_{0}
=\bigl(\mathcal{L}\mathrm{ie}(X_{0}/S_{0})
/\mathcal{L}\mathrm{ie}(Y_{0}/S_{0})\bigr)
\otimes_{\OO_{S_{0}}}\mathcal{J}\),
one has \(E=\Gamma(S_{0},\mathcal{F}_{0})/H^{0}(Y_{0},\mathcal{F}_{0})\).
\end{corollary}

{\renewcommand{\thefootnote}{(128)}\nde{One has spelled out the original in what follows.}}
Indeed, since \(\omega^{1}_{Y_{0}/S_{0}}\) is finite locally free, as is
\(\omega^{1}_{X_{0}/S_{0}}\) (since \(X\) is assumed smooth over~\(S\)), one
has, by~0.6:
\[
L_{0Y}
=
\mathbf{W}\bigl(\mathcal{L}\mathrm{ie}(Y_{0}/S_{0})\otimes_{\OO_{S_{0}}}\mathcal{J}\bigr)
\qquad\text{and}\qquad
L_{0X}
=
\mathbf{W}\bigl(\mathcal{L}\mathrm{ie}(X_{0}/S_{0})\otimes_{\OO_{S_{0}}}\mathcal{J}\bigr)\,.
\]
On the other hand, by~4.25, one has an exact sequence of
\(Y_{0}\)-\(\OO_{S_{0}}\)-modules:
\[
0\longrightarrow\mathcal{K}\longrightarrow
\omega^{1}_{X_{0}/S_{0}}
\xrightarrow{\phi}
\omega^{1}_{Y_{0}/S_{0}}
\longrightarrow 0
\]
(where \(\mathcal{K}=\Ker(\phi)\)). Since \(\omega^{1}_{Y_{0}/S_{0}}\) and
\(\omega^{1}_{X_{0}/S_{0}}\) are finite locally free, one has a locally split
exact sequence:
\[
0\longrightarrow
\mathcal{L}\mathrm{ie}(Y_{0}/S_{0})\otimes_{\OO_{S_{0}}}\mathcal{J}
\longrightarrow
\mathcal{L}\mathrm{ie}(X_{0}/S_{0})\otimes_{\OO_{S_{0}}}\mathcal{J}
\longrightarrow
\mathcal{F}_{0}
\longrightarrow 0
\]
It follows that one has an exact sequence of
\(Y_{0}\)-\(\OO_{S_{0}}\)-modules:
\[
0\longrightarrow L_{0Y}\longrightarrow L_{0X}\longrightarrow\mathbf{W}(\mathcal{F}_{0})\,.
\]
By the argument which we used to prove~4.31, we may compute \(E\) as the
image of the composite map
\[
C^{0}(Y_{0},L_{0X})
\xrightarrow{\pi}
C^{0}(Y_{0},\,\mathbf{W}(\mathcal{F}_{0}))
\xrightarrow{\partial}
C^{1}(Y_{0},\,\mathbf{W}(\mathcal{F}_{0}))\,.
\]
Now the map \(\pi\) above is the map
\(\Gamma\bigl(S_{0},\,
\mathcal{L}\mathrm{ie}(X_{0}/S_{0})\otimes_{\OO_{S_{0}}}\mathcal{J}\bigr)
\to\Gamma(S_{0},\mathcal{F}_{0})\).
Thus, \(S_{0}\) being affine, \(\pi\) is surjective, and one indeed finds the
result announced.

\begin{corollary}{4.33}\label{III.4.33}
Let \(S\), \(\mathcal{I}\), \(\mathcal{J}\) and \(X\) be as in~4.21, and let
\(Y\) be a diagonalizable subgroup of~\(X\). Suppose that
\(\omega^{1}_{Y_{0}/S_{0}}\) is finite locally free and that \(S_{0}\) is
affine.%
{\renewcommand{\thefootnote}{(129)}\nde{\normalfont One has added the hypothesis on \(\omega^{1}_{Y_{0}/S_{0}}\) and replaced the hypothesis ``\(S\) affine'' by ``\(S_{0}\) affine''.}}
The set of subgroups of~\(X\), conjugates of \(Y\) by a section of \(X\)
over \(S\) inducing the unit section of \(X_{\mathcal{J}}\), is isomorphic to
\[
E
=
\Gamma\Bigl(S_{0},\,
\mathcal{L}\mathrm{ie}(X_{0}/S_{0})
\big/
\mathcal{L}\mathrm{ie}(X_{0}/S_{0})^{\mathrm{ad}(Y_{0})}
\Bigr)
\otimes_{\Gamma(S_{0},\OO_{S_{0}})}
\Gamma(S_{0},\mathcal{J})
\]
{\renewcommand{\thefootnote}{(130)}\nde{\normalfont One has added what follows, \cf~4.34.}}
that is, to \(L_{0X}(Y_{0})/H^{0}(Y_{0},L_{0X})\).
\end{corollary}

Indeed, one writes, by~I 4.7.3 (\cf~2.13):
\[
\mathcal{L}\mathrm{ie}(X_{0}/S_{0})
=
\mathcal{L}\mathrm{ie}(X_{0}/S_{0})^{\mathrm{ad}(Y_{0})}
\oplus
\mathcal{R}\,.
\]
Since \(Y_{0}\) is commutative, one has
\(\mathcal{L}\mathrm{ie}(Y_{0}/S_{0})
\subset
\mathcal{L}\mathrm{ie}(X_{0}/S_{0})^{\mathrm{ad}(Y_{0})}\),
hence
\begin{align*}
\mathcal{F}_{0}
&=
\bigl(
\mathcal{L}\mathrm{ie}(X_{0}/S_{0})^{\mathrm{ad}(Y_{0})}
\big/
\mathcal{L}\mathrm{ie}(Y_{0}/S_{0})
\bigr)
\otimes\mathcal{J}
\oplus
\mathcal{R}\otimes\mathcal{J}\,,
\\
\mathcal{F}_{0}^{\mathrm{ad}(Y_{0})}
&=
\bigl(
\mathcal{L}\mathrm{ie}(X_{0}/S_{0})^{\mathrm{ad}(Y_{0})}
\big/
\mathcal{L}\mathrm{ie}(Y_{0}/S_{0})
\bigr)
\otimes\mathcal{J}\,.
\end{align*}
By~4.32, one therefore has \(E\simeq\Gamma(S_{0},\,\mathcal{R}\otimes\mathcal{J})\).
Returning to the definition of \(\mathcal{R}\), one is done.

\begin{corollary}{4.34}\label{III.4.34}
Let \(S\), \(\mathcal{I}\), \(\mathcal{J}\) and \(X\) be as in~4.21, and let
\(Y\) be a diagonalizable subgroup of~\(X\). Suppose that
\(\omega^{1}_{Y_{0}/S_{0}}\) is finite locally free and that \(S_{0}\) is
affine.%
{\renewcommand{\thefootnote}{(131)}\nde{\normalfont \Cf\ N.D.E.~(129).}}
If \(x\in X(S)\) induces the unit section of \(X_{\mathcal{J}}\) and
normalizes~\(Y\), then it centralizes~\(Y\).
\end{corollary}

This follows immediately from the comparison of the preceding corollary and
of~2.14. Indeed, 4.33 shows that the elements of \(C^{0}(Y_{0},L_{0X})\)
which globally preserve \(Y\) are the elements of \(H^{0}(Y_{0},L_{0X})\), and
one has seen in~2.14 that they are precisely those which act trivially on
the canonical immersion \(Y\to X\).

\numpara{4.35}\label{III.4.35}
Let us return to the general situation of~4.21 and assume \(Y_{\mathcal{J}}\)
smooth over \(S_{\mathcal{J}}\). Then, by~4.25~(ii), one has an exact
sequence of \(Y_{0}\)-\(\OO_{S_{0}}\)-modules:
\[
(*)\qquad
0\longrightarrow
\mathcal{L}\mathrm{ie}(Y_{0}/S_{0})
\longrightarrow
\mathcal{L}\mathrm{ie}(X_{0}/S_{0})
\longrightarrow
\mathfrak{n}^{\vee}_{Y_{0}/X_{0}}
\longrightarrow 0
\]
and these are finite locally free \(\OO_{S_{0}}\)-modules.

On the other hand, by SGA~1, II~4.10, every subscheme \(Y\) of \(X\) lifting
\(Y_{\mathcal{J}}\) and \emph{flat} over \(S\) will be \emph{smooth} over~\(S\).%
{\renewcommand{\thefootnote}{(132)}\nde{One has added what follows and Proposition~4.35.1, implicit in the original, \cf~4.38~(5).}}
Suppose moreover that \(S_{0}\) and \(Y_{\mathcal{J}}\) are \emph{affine}.
Then, since \(\mathfrak{n}^{\vee}_{Y_{0}/X_{0}}\) is a locally free
\(\OO_{S_{0}}\)-module, the sequence \((*)\) remains exact when one applies
\(\otimes_{\OO_{S_{0}}}\mathcal{J}\) to it, and then takes the inverse image
on \(Y_{0}^{n}\), and since the \(Y_{0}^{n}\) are affine, one therefore
obtains an exact sequence of complexes of abelian groups:
\[
0\longrightarrow
C^{*}(Y_{0},L_{0Y})\longrightarrow
C^{*}(Y_{0},L_{0X})\xrightarrow{d^{*}}
C^{*}(Y_{0},N_{0})\longrightarrow 0
\]
and in particular, a commutative diagram with exact rows
\[
\xymatrix@C=1.45em@R=1.8em{
0\ar[r]
&
C^{1}(Y_{0},L_{0Y})
  \ar[r]
  \ar[d]^{\partial}
&
C^{1}(Y_{0},L_{0X})
  \ar[r]^{d^{1}}
  \ar[d]^{\partial}
&
C^{1}(Y_{0},N_{0})
  \ar[r]
  \ar[d]^{\partial}
&
0
\\
0\ar[r]
&
C^{2}(Y_{0},L_{0Y})
  \ar[r]
&
C^{2}(Y_{0},L_{0X})
  \ar[r]^{d^{2}}
&
C^{2}(Y_{0},N_{0})
  \ar[r]
&
0\,.
}
\]
Now let \(Y\), \(Y'\) be two subgroups of \(X\) lifting \(Y_{\mathcal{J}}\) and
\emph{flat}, hence \emph{smooth}, over~\(S\). Since \(Y_{\mathcal{J}}\) is
affine, then, by~0.15, \(Y\) and \(Y'\) are isomorphic as schemes extending
\(Y_{\mathcal{J}}\), \ie\ there exists an isomorphism of \(S\)-schemes
\[
f\colon Y\isomto Y'
\]
inducing the identity on \(Y_{\mathcal{J}}\).
